\section{Related Work}
\label{sec:related}

\subsection{Multi-Dimensional Trustworthiness Evaluation Frameworks}

Modern LLM evaluation has progressed from single-task accuracy benchmarks toward comprehensive multi-dimensional assessment. TrustLLM \cite{sun2024trustllm} provides the most complete current framework, evaluating 16 prominent LLMs across 6 trustworthiness dimensions (truthfulness, safety, fairness, adversarial robustness, privacy, machine ethics) using a suite of sub-benchmarks. HELM \cite{liang2023helm} offers a complementary perspective with 7 metrics (accuracy, calibration, robustness, fairness, bias, toxicity, efficiency) across 30 models. MultiTrust \cite{zhang2024multitrust} extends the analysis to multimodal LLMs across 5 dimensions. Crucially, all three frameworks present dimension scores and visual radar charts but do not compute the pairwise correlation structure between dimensions or test whether dimensions are statistically independent across models.

Liu et al.\ \cite{liu2023trustworthy} survey trustworthiness across 7 categories and explicitly identify cross-dimension correlation analysis as future work. Our paper delivers that analysis. In this sense, we do not compete with these evaluation frameworks --- we analyze the data they provide.

\subsection{Trustworthiness Tradeoffs: Specific Pairs}

Several studies have identified specific trustworthiness tradeoffs without characterizing the full correlation structure. AQUA-LLM \cite{gungor2025aquallm} demonstrates a quantitative accuracy-robustness tradeoff under quantization and adversarial perturbation --- corroborating the directional safety-robustness tension we observe. Know Thy Judge \cite{eiras2025know} shows that RLHF-trained safety judges are brittle to style shifts, with false negative rates jumping by up to 0.24 --- consistent with our finding that RLHF-optimized models may not transfer their safety behaviors to adversarial inputs. Wang et al.\ \cite{wang2025survey} survey reasoning LLMs and find that chain-of-thought models show \emph{worse} safety and privacy performance than standard models at equivalent scale, suggesting inverse scaling relationships for specific dimensions.

These studies confirm that trustworthiness dimensions are not uniformly correlated --- some pairs trade off while others co-move. However, they focus on specific pairs rather than the full 15-pair correlation matrix, and none controls for the scale and RLHF confounds that conflate model quality with trustworthiness geometry.

\subsection{Benchmark Correlation Analysis and Evaluation Compression}

The question of whether benchmark scores are redundant is well-studied for general capability benchmarks. The Epoch AI benchmark correlation study \cite{epochai2024} finds a median Spearman $\rho = 0.73$ across 17 capability benchmarks, suggesting substantial redundancy in capability measurement.

Minimum spanning tree methods for correlation-based redundancy analysis have been applied to financial return correlation matrices \cite{tumminello2005tool} and ecological networks, where bootstrap stability metrics characterize the reliability of inferred tree topology. We apply this methodology to the LLM trustworthiness domain for the first time, using the Tumminello et al.\ (2005) mean per-edge bootstrap frequency criterion to characterize MST stability.

\textbf{Note on the Epoch AI baseline:} We use the Epoch AI capability baseline (median $\rho = 0.73$) as a reference point for the scale of benchmark redundancy in the field. However, the Epoch AI study reports \emph{raw} Spearman correlations among capability benchmarks, while the correlations we report are \emph{partial} Spearman after confound removal. These two quantities are methodologically distinct: partial correlations structurally differ from raw correlations when confounds are present (the paper demonstrates this for the safety--privacy pair, where raw $\rho = 0.73$ vs.\ partial $\rho = 0.971$). Accordingly, the Epoch AI baseline should be read as a reference point for the scale of correlation in the broader benchmark literature, not as a methodologically equivalent comparison.

\subsection{RLHF Effects on Trustworthiness}

The RLHF alignment paradigm is the primary mechanism through which modern LLMs are adapted for safe and helpful behavior \cite{ouyang2022training}. Li et al.\ \cite{li2025more} (ICLR Oral) find that ``more RLHF'' does not automatically guarantee trustworthiness across all dimensions --- safety can improve while other dimensions stagnate or regress. Our findings are consistent with Li et al.: we confirm that RLHF jointly optimizes safety and ethics (via LLaMA-2 within-family evidence) while robustness follows an independent path.

\subsection{Positioning Our Contribution}

\begin{table}[h]
\centering
\caption{Gap analysis: what prior work provides vs.\ what is missing.}
\label{tab:gap}
\begin{tabular}{lll}
\toprule
\textbf{Prior Work} & \textbf{Data Provided} & \textbf{What is Missing} \\
\midrule
TrustLLM \cite{sun2024trustllm} & 16-model $\times$ 6-dim scores & Pairwise correlation; no partial corr. \\
HELM \cite{liang2023helm} & 30-model $\times$ 7-metric scores & Cross-metric matrix; no clustering \\
Liu et al.\ \cite{liu2023trustworthy} & 7-category survey & Lists cross-dim corr.\ as future work \\
AQUA-LLM \cite{gungor2025aquallm} & Accuracy-robustness tradeoff & Single pair; no full matrix \\
Epoch AI \cite{epochai2024} & Capability benchmark corr.\ (raw $\rho$=0.73) & Capability, not trustworthiness \\
\bottomrule
\end{tabular}
\end{table}

Our work fills the gap at the intersection: the first pairwise partial Spearman correlation matrix for LLM trustworthiness, the first cluster analysis of the correlation geometry, and the first MST-based minimum evaluation set --- all using publicly available TrustLLM data.
